\scanpage{165}
% L1–9
Cherchons les formes quadratiques \unsure{(E}
sur $E$ compatibles avec une involution déjà donnée, \inserted{i.e. avec un $t\in E$ donné.}
Il est immédiat \unsure{(.|(}
que ces formes forment un torseur
sous l'\unsure{ens|un} des formes \inserted{(à val. dans $L^{\otimes 2}$)}
qui se factorisent par $E/L\simeq k$, i.e. sous
$\deleted{\unsure{k}}L^{\otimes 2}$ lui $=$.
On en a déjà un tel torseur, savoir $A*A$, et une
application quadratique \unsure{canonique|homogène}
\[
A\deleted{*A}\longrightarrow A*A
\]
% L11–14
(``compatible avec chg$^{\mathrm t}$ de base'' donc
\unsure{défini|defini|définie} sur les schémas correspondants)
savoir
\[
x\,\unsure{\mathrel{\vdash}|\mapsto}\,x*x
\]

% L15–22
Ainsi, \deleted{\unsure{soit}} pour \deleted{tout}
\[
\zeta\in A*A,
\]
si on dénote par $q_\zeta$ l'unique \deleted{forme} \inserted{fonction}
\deleted{quadratique} \inserted{polyn. de degré $\leq 2$}
\[
q_\zeta:\deleted{\unsure{E}}A\longrightarrow L^{\otimes 2}
\]
qui satisfait
\[
\zeta\deleted{\unsure{-}}+q_\zeta(x)=x*x\qquad \forall\,x\text{ dans }A
\]
on aura, pour $\alpha\in L^{\otimes 2}$, \unsure{remplaçant}
$\zeta$ par $\zeta+\alpha$
\[
(\zeta+\alpha)+q_{\zeta+\alpha}(\unsure{u|x})=x*x
\]
donc
\[
q_{\zeta+\alpha}(x)\deleted{\unsure{+}}-q_\zeta(x)=
\underset{\substack{\Vert\\-\alpha\epsilon(x)^2}}{\deleted{\unsure{+}}-\alpha}
\quad\begin{aligned}&\text{pour }x\in A\\&\text{i.e. }\epsilon(x)=1\end{aligned}
\]
% archive sheet number: 164
